\section*{Capítulo \ref{ch:b1:diagenesis-basins} --- Diagênese e bacias sedimentares}

\begin{solution}{exo:b1:diagenesis-basins:1}
Arenito $0.49\,e^{-2000/3704} = 0.49\times0.583 = \num{0.286}$; folhelho
$0.63\,e^{-2000/1960} = 0.63\times0.360 = \num{0.227}$. A porosidade cai à metade quando
$e^{-z/\lambda} = 1/2$, em $z = \lambda\ln 2$: \qty{2567}{m} para o arenito,
\qty{1359}{m} para o folhelho.
\end{solution}

\begin{solution}{exo:b1:diagenesis-basins:2}
$\Delta A = \Delta E + S = 30 + 50 = \qty{80}{m}$. Dele, \qty{60}{m} são preenchidos com
sedimento, então $\Delta W = 80 - 60 = \qty{20}{m}$: a plataforma aprofunda-se em
\qty{20}{m}.
\end{solution}

\begin{solution}{exo:b1:diagenesis-basins:3}
$h = Y(\rho_m - \rho_w)/(\rho_m - \rho_s)$: com \qty{2300}{kg.m^{-3}},
$1.0\times2270/1000 = \qty{2.27}{km}$; com \qty{2500}{kg.m^{-3}},
$1.0\times2270/800 = \qty{2.84}{km}$. Cada metro de sedimento ocupa o lugar de um metro de
\omterm{def:b1:layered-earth:layers}{manto} na coluna; quanto mais denso o sedimento, menor a massa que lhe falta em relação ao
\omterm{def:b1:layered-earth:layers}{manto}, de modo que mais dele é necessário para compensar o mesmo déficit que o buraco
cheio de água.
\end{solution}

\begin{solution}{exo:b1:diagenesis-basins:4}
(a)~Calcita $100.09/2.71 = \qty{36.93}{cm^3.mol^{-1}}$; dolomita $184.40/2.84
= \qty{64.93}{cm^3.mol^{-1}}$. (b)~Dois mols de calcita ocupam
\qty{73.86}{cm^3}, um de dolomita \qty{64.93}{cm^3}: o sólido encolhe
$1 - 64.93/73.86 = \qty{12.1}{\%}$. (c)~Em um metro cúbico, os \qty{0.90}{m^3}
de calcita tornam-se $0.90\times(1 - 0.121) = \qty{0.791}{m^3}$ de dolomita; a
porosidade torna-se $1 - 0.791 = \num{0.21}$.
\end{solution}

\begin{solution}{exo:b1:diagenesis-basins:5}
(a)~$s = 1000 - 0.63\times1960\times(1 - e^{-1000/1960}) = \qty{506.5}{m}$; a
porosidade média é $1 - 506.5/1000 = \num{0.493}$. (b)~Com o topo a
\qty{2000}{m}, resolver $s(2000, b) = 506.5$ dá $b = \qty{2628.7}{m}$: a
camada tem \qty{628.7}{m} de espessura, com porosidade média $1 - 506.5/628.7 = \num{0.194}$.
(c)~Os grãos são os mesmos, então a água perdida é a espessura perdida:
$1000 - 628.7 = \qty{371}{m^3}$ de água por metro quadrado.
\end{solution}

\begin{solution}{exo:b1:diagenesis-basins:6}
(a)~$s = 300 - 0.49\times3704\times(e^{-2500/3704} - e^{-2800/3704}) =
\qty{228.10}{m}$. (b)~A partir de $b_0 = 300$: $s(0, 300) = 158.80$, $1 - \varphi(300) =
0.5481$, $b_1 = 300 + (228.10 - 158.80)/0.5481 = \qty{426.44}{m}$; então
$s(0, 426.44) = 229.07$, $1 - \varphi = 0.5633$, $b_2 = 426.44 - 0.96/0.5633 =
\qty{424.73}{m}$; um terceiro passo muda isso em menos de um milímetro. O arenito tinha
\qty{424.7}{m} de espessura na superfície. (c)~$424.7/300 = 1.42$: o soterramento o
afinou por um fator de \num{1.42}.
\end{solution}

\begin{solution}{exo:b1:diagenesis-basins:7}
(a)~O gelo contém $52\times10^6\times916.7/1000 = \qty{4.77e7}{km^3}$ de água;
espalhado sobre \qty{361.84e6}{km^2} isso dá $\qty{0.1317}{km} = \qty{132}{m}$ de nível
do mar. (b)~As linhas de costa dão \qty{134}{m}, \qty{2}{m} a mais. (c)~$4.767\times
10^7/0.134 = \qty{355.7e6}{km^2}$, \qty{1.7}{\%} a menos que hoje. À medida que o nível
do mar caía, as plataformas continentais emergiam, de modo que o oceano glacial era menor,
e o mesmo volume de água retirado dele o rebaixava mais.
\end{solution}

\begin{solution}{exo:b1:diagenesis-basins:8}
(a)~$(-9.1 - (-98.1))/8 = \qty{11.1}{m}$ por mil anos, perto dos \qty{12}{m} encontrados
a partir das linhas de costa. (b)~\qty{40}{mm} por ano são \qty{40}{m} por mil anos,
$40/12 = 3.3$ vezes a média; em quinhentos anos, $0.040\times500 = \qty{20}{m}$.
\end{solution}

\begin{solution}{exo:b1:diagenesis-basins:9}
(a)~$S = 50\times2 = \qty{100}{m}$ e $\Delta E = \qty{-40}{m}$:
$\Delta A = \qty{60}{m}$, todo preenchido com sedimento, já que a profundidade não muda.
(b)~Em \num{20000} anos o embasamento afunda apenas $50\times0.02 = \qty{1}{m}$:
o \omterm{def:b1:diagenesis-basins:eustasy}{nível relativo do mar} cai $120 - 1 = \qty{119}{m}$. O fundo marinho, que estava
\qty{20}{m} abaixo do nível do mar, termina $119 - 20 = \qty{99}{m}$ acima dele: a
plataforma inteira emerge.
\end{solution}

\begin{solution}{exo:b1:diagenesis-basins:10}
(a)~Para uma camada fina a porosidade é quase constante através dela: sua espessura sólida é
$s \approx h(1 - \varphi(z))$ na profundidade e $s \approx h_0(1 -
\varphi_0)$ na superfície. Igualando as duas obtém-se $h_0 \approx h(1 -
\varphi(z))/(1 - \varphi_0)$. (b)~$\varphi(3000) = 0.63\,e^{-3000/1960} =
0.136$, logo $h_0 \approx 10\times0.864/0.37 = \qty{23.3}{m}$, contra
\qty{23.1}{m} exatos. (c)~Restaurada à superfície, a camada tem \qty{23}{m} de espessura,
já não é fina em comparação com o decaimento da porosidade perto da superfície: sua
porosidade cai de $\varphi_0$ no topo para menos na base, de modo que seus grãos cabem em
um pouco menos que a espessura calculada com $\varphi_0$ por toda parte.
\end{solution}

\begin{solution}{exo:b1:diagenesis-basins:11}
(a)~$s = 3000 - 0.63\times1960\times(1 - e^{-3000/1960}) = \qty{2032}{m}$; a porosidade
média é $1 - 2032/3000 = \num{0.323}$ e a densidade média
$0.323\times1030 + 0.677\times2700 = \qty{2161}{kg.m^{-3}}$. (b)~O fator é
$(3300 - 1030)/(3300 - 2161) = \num{1.99}$; o preenchimento representa uma \omterm{def:b1:diagenesis-basins:subsidence}{subsidência} com
água de $3000/1.99 = \qty{1505}{m}$. (c)~Com a densidade dos grãos o fator seria
$2270/600 = \num{3.78}$ e a \omterm{def:b1:diagenesis-basins:subsidence}{subsidência} apenas \qty{793}{m}. Isso estaria errado porque um
terço do preenchimento é água: o sedimento é muito mais leve que seus grãos, e só a
densidade do sedimento saturado conta no balanço de massas.
\end{solution}

\begin{solution}{exo:b1:diagenesis-basins:12}
(a)~O cubo contém \qty{0.90}{m^3} $= \qty{9.0e5}{cm^3}$ de calcita, isto é
$9.0\times10^5/36.93 = \num{24370}$ mols. (b)~Um magnésio para cada duas calcitas:
\num{12180} mols. (c)~Um quilograma de água do mar contém $35.165\times0.0365 =
\qty{1.284}{g}$ de magnésio, $1.284/24.305 = \qty{0.0528}{mol}$; a reação precisa de
$12180/0.0528 = \qty{2.31e5}{kg}$ de água do mar, $\qty{224}{m^3}$ --- mais de duzentas
vezes o volume da rocha, mesmo que todo íon de magnésio seja tomado. (d)~A dolomita só pode
se formar em grande escala onde volumes muito grandes de água fluem pela rocha por muito
tempo, ou onde a água é uma salmoura mais rica em magnésio que a água do mar: em plataformas
inundadas por salmouras que afundam, por exemplo, e não em um sedimento que mantém a mesma
água de poros.
\end{solution}

\begin{solution}{pb:b1:diagenesis-basins:1}
\textbf{1.} Arenito: $\varphi(0) = \num{0.490}$, $\varphi(800) =
0.49\,e^{-800/3704} = \num{0.395}$. Folhelho: $\varphi(800) = \num{0.419}$,
$\varphi(2600) = \num{0.167}$. Giz: $\varphi(2600) = \num{0.110}$,
$\varphi(3200) = \num{0.072}$.

\textbf{2.} Com $s$ de \cref{prop:b1:diagenesis-basins:solid} e $\bar\varphi
= 1 - s/h$ (questão~7): arenito \num{0.441}, folhelho \num{0.274}, giz \num{0.090}.

\textbf{3.} $h\bar\varphi$: arenito \qty{353}{m^3}, folhelho \qty{493}{m^3}, giz
\qty{54}{m^3} de água por metro quadrado.

\textbf{4.} $\bar\varphi\rho_w + (1 - \bar\varphi)\rho_g$: arenito
$0.441\times1030 + 0.559\times2650 = \qty{1936}{kg.m^{-3}}$, folhelho
\qty{2242}{kg.m^{-3}}, giz \qty{2559}{kg.m^{-3}}. A coluna pesa
$1936\times800 + 2242\times1800 + 2559\times600 = \qty{7.12e6}{kg}$ por metro
quadrado: uma densidade média de \qty{2225}{kg.m^{-3}}.

\textbf{5.} Litostática $2225\times9.81\times3200 = \qty{69.9}{MPa}$; hidrostática
$1030\times9.81\times3200 = \qty{32.3}{MPa}$; \omterm{def:b1:stress-strain-rheology:lithostatic}{tensão efetiva}
$69.9 - 32.3 = \qty{37.5}{MPa}$.

\textbf{6.} $0.49\,e^{-z/3704} = 0.63\,e^{-z/1960}$ dá $z = \ln(0.63/0.49)
/(1/1960 - 1/3704) = \qty{1046}{m}$. Abaixo dela o arenito é o mais poroso: o folhelho
começa mais úmido, mas perde sua água em uma distância menor.

\textbf{7.} Arenito $800 - 0.49\times3704\times(1 - e^{-800/3704}) =
\qty{447.4}{m}$; folhelho \qty{1306.7}{m}; giz \qty{546.0}{m}.

\textbf{8.} \qty{2300.2}{m} de grãos; a água é $1 - 2300.2/3200 =
\qty{28}{\%}$ da coluna.

\textbf{9.} Arenito $447.4/15 = \qty{29.8}{m}$, folhelho $1306.7/40 =
\qty{32.7}{m}$, giz $546.0/20 = \qty{27.3}{m}$ de sólido por milhão de anos.

\textbf{10.} O folhelho teve o suprimento mais rápido, o giz o mais lento --- mas os três
diferem em apenas cerca de um quinto, muito menos do que suas espessuras hoje sugerem.

\textbf{11.} \qty{800}{m}: seu topo está na superfície hoje, onde sua deposição acabou de
terminar, de modo que sua espessura na deposição é a espessura de hoje.

\textbf{12.} Resolvendo $s(0, b) = 1306.7$ para o folhelho por
\cref{met:b1:diagenesis-basins:decompaction} obtém-se $b = \qty{2124}{m}$: o folhelho
tinha \qty{2124}{m} de espessura a \qty{15}{\Ma}.

\textbf{13.} $s(0, b) = 546.0$ para o giz dá \qty{1071}{m} a \qty{55}{\Ma}.

\textbf{14.} Com seu topo a \qty{2124}{m}, $s(2124, b) = 546.0$ dá $b =
\qty{2748}{m}$: o giz tinha \qty{624}{m} de espessura.

\textbf{15.} \qty{1071}{m} a \qty{55}{\Ma}; $2124 + 624 = \qty{2748}{m}$ a \qty{15}{\Ma}.

\textbf{16.} Pelas espessuras de hoje: arenito $800/15 = \qty{53.3}{m}$, folhelho
$1800/40 = \qty{45.0}{m}$, giz $600/20 = \qty{30.0}{m}$ por milhão de anos. Pelas
espessuras na deposição: \qty{53.3}{m}, $2124/40 =
\qty{53.1}{m}$ e $1071/20 = \qty{53.6}{m}$. O segundo conjunto é o significativo --- a
taxa com que o fundo do mar foi construído --- e mostra uma taxa constante que as
espessuras de hoje escondem.

\textbf{17.} $1 - 1800/2124 = \qty{15}{\%}$.

\textbf{18.} A profundidade da água é a mesma no início e no fim de cada intervalo, então a
\omterm{def:b1:diagenesis-basins:accommodation}{acomodação} criada é o crescimento da coluna acima do datum: \qty{1071}{m} de 75 a
\qty{55}{\Ma}, $2748 - 1071 = \qty{1677}{m}$ de 55 a \qty{15}{\Ma}, e
$3200 - 2748 = \qty{452}{m}$ desde \qty{15}{\Ma}.

\textbf{19.} $S = \Delta A - \Delta E$: $1677 - (-60) = \qty{1737}{m}$ de 55 a
\qty{15}{\Ma}; $452 - (-30) = \qty{482}{m}$ desde então.

\textbf{20.} $1737/40 = \qty{43}{m}$ e $482/15 = \qty{32}{m}$ por milhão de anos.

\textbf{21.} $60/1737 = \qty{3.5}{\%}$ e $30/482 = \qty{6.2}{\%}$. O nível do mar caiu nos
dois intervalos, então a \omterm{def:b1:diagenesis-basins:eustasy}{eustasia} atrapalhou o preenchimento: o embasamento teve de afundar
essa quantidade a mais para criar o mesmo espaço. A \omterm{def:b1:diagenesis-basins:subsidence}{subsidência} criou essencialmente toda a
\omterm{def:b1:diagenesis-basins:accommodation}{acomodação}.

\textbf{22.} O folhelho e o giz tinham \qty{2748}{m} de espessura a \qty{15}{\Ma} e hoje têm
\qty{2400}{m}: sua \omterm{def:b1:diagenesis-basins:compaction}{compactação} criou \qty{348}{m} de espaço, \qty{77}{\%} dos \qty{452}{m}
preenchidos pelo arenito.

\textbf{23.} Cada unidade foi afinada pelo soterramento que se seguiu, por uma quantidade que
depende de sua litologia e de quão fundo ela foi; o espaço criado em cada intervalo deve ser
medido com as espessuras que as unidades tinham então, o que só a \omterm{def:b1:diagenesis-basins:compaction}{descompactação} fornece. E
mesmo assim a \omterm{def:b1:diagenesis-basins:subsidence}{subsidência} assim encontrada inclui o afundamento causado pelo peso do
sedimento, que deve ser removido para chegar à \omterm{def:b1:diagenesis-basins:subsidence}{subsidência tectônica}.

\textbf{24.} \qty{2124}{m}, cerca de \qty{2.1}{km} --- \qty{324}{m} a mais que os
\qty{1800}{m} que o poço encontra hoje.
\end{solution}
